\chapter*{Preface}\label{ch-preface}
\addcontentsline{toc}{chapter}{Preface}
\markboth{Preface}{Preface}

Zeno is a programming language that reads like plain English and is
checked like Rust. It was built from nothing, one stage at a time, by
a firmware engineer who wanted to know how a language works from the
inside: the lexer, the parser, the type checker, two engines that run
programs, two compilers that turn them into machine code, and the
library and tools around them. This book is the whole of that
language, written down once it stopped moving.

\subsection*{Who this book is for}

The book starts from nothing. If you have never written a program,
start at the first page and read in order: every chapter uses only
what the chapters before it showed you.

If you already know another language, you will find Zeno familiar
quickly. It takes its spelling from Python wherever Python's spelling
is clear, its braces from C, its checking from Rust, its sized numbers
from firmware, and its error messages from the idea that a message
should say what to write instead. Read Part~I quickly, skim Part~II
for the places where Zeno differs from what you know, and slow down
from Part~III on.

If you write firmware, Part~IV is for you: it shows how one Zeno
program builds for a computer, for a firmware coding standard with no
\verb|setjmp|, with no garbage collector at all, and with memory
handled by hand, and how each of those builds is told apart.

If you want to know how the compiler itself works, Part~V walks
through it, file by file.

\subsection*{How the book is laid out}

\begin{description}
\item[Part I, Starting] gets Zeno onto your machine, runs a first
  program, and takes a quick tour of the whole language so you know
  what is coming.
\item[Part II, The Language] is everything a program is made of:
  values, numbers, text, decisions and loops, functions, lists,
  dicts, classes, enums, values that may be missing, errors, and
  functions used as values.
\item[Part III, Growing Programs] is what a program needs once it is
  more than a page: types written and worked out, units, sized
  numbers and bits, generics, interfaces, modules, packages, files,
  mathematics, tests, and the tools.
\item[Part IV, Zeno on a Machine] is about where a program runs: the
  board and its pins, compiling to C and to LLVM, memory and the
  collector, owned values, firmware builds, builds with no collector,
  memory by hand, and a program in a web page.
\item[Part V, Inside Zeno] is how the compiler is built and how the
  language is designed.
\item[The appendices] are for looking things up: every keyword, every
  operator, the whole library, every unit, every command, every kind
  of error, and the list of what Zeno does not have yet.
\end{description}

Each chapter ends with \emph{Things to try}: small changes to make and
questions to answer by running Zeno. They are the fastest way to make
a chapter your own.

\subsection*{The programs in this book are real}

Every program in this book was run by the Zeno compiler when the book
was built, and every output under a program is what Zeno printed. That
includes the error messages: when the book shows a mistake, the
message under it is the one Zeno gives. If a later version of Zeno
prints something else, the book is rebuilt and the page changes with
it.

A program looks like this:

\begin{zeno}
print("Hello, Zeno!")
\end{zeno}

and what it printed is in the paler box under it. When a program is
written to a file of its own name, the name is shown above it. A
command typed into a terminal starts with \verb|$|, and its output
follows on the next lines.

\subsection*{The boxes}

\begin{note}
A note says something worth knowing that the text around it does not
stop to explain.
\end{note}

\begin{why}
Zeno has a recorded reason for every choice in it. A box like this
gives the reason for the rule beside it, and often what was turned
down. All of the reasons are kept in \verb|docs/decisions.md| in the
Zeno repository, with a number for each, and the book names the number
when it is useful.
\end{why}

\subsection*{When you are stuck}

If a step in this book does not work for you, or you cannot see how to
write something, write to \texttt{help@zeno-lang.org}. Say what you
typed and what Zeno printed.

If you think you have found a mistake in Zeno itself, report it as an
issue at \url{https://github.com/partham20/zeno-lang/issues}, where
others can see it and its answer. If you have no GitHub account, write
to \texttt{issues@zeno-lang.org} instead. A mistake in this book, an
idea, or anything else goes to \texttt{contact@zeno-lang.org}.

\subsection*{The rule above the rules}

One sentence decides more about Zeno than any other: \emph{the entire
language must be unambiguous and intuitive}. In practice that means
five things, and you will see each of them again and again in this
book:

\begin{enumerate}
\item There is one obvious way to write each thing.
\item No line can be read two ways, by the compiler or by a person.
\item Nothing happens that the line does not say.
\item Every word is a plain word, with one meaning.
\item When Zeno refuses something, the message says what to write instead.
\end{enumerate}

\subsection*{The name}

Zeno of Elea argued that Achilles can never catch the tortoise: before
he reaches it he must cover half the distance, then half of what is
left, and so on forever. The quantum Zeno effect says that a system
watched closely enough never changes. Both suit a compiler that will
always be almost finished. Its motto is the first thing \verb|zeno --help| prints:

\begin{quote}
\emph{Every journey is halfway done. Forever.}
\end{quote}
